\chapter{Меню spectral functions и moment-selector gate}

Compact \(a_2/a_4\) calculation свёл scale-setting к комбинации
\begin{equation}
 \zeta_{\mathrm{mom}}
 =
 \frac{f_2}{f_0}(\Lambda R)^2,
 \qquad
 (\chi R)^2=\zeta_{\mathrm{mom}}-\frac12.
\end{equation}
Теперь проверим, фиксируют ли стандартные требования к spectral function
отношение \(f_2/f_0\).

\section{Moment convention}

Пусть spectral action записан как
\begin{equation}
 \Tr f(\mathcal D^2/\Lambda^2),
\end{equation}
где \(u\ge0\) --- безразмерный аргумент. Используем
\begin{equation}
 f_0=f(0),
 \qquad
 f_2=\int_0^\infty f(u)\,du.
\end{equation}
Общая scalar normalization позволяет положить
\begin{equation}
 f_0=1
\end{equation}
без потери relative coefficients. Тогда relevant datum равно площади под
профилем \(f\).

\section{Пререгистрированное минимальное меню}

До обращения к particle masses фиксируем четыре стандартных профиля:
\begin{align}
 f_{\mathrm{sharp}}(u)
 &=\mathbf1_{[0,1]}(u),
 \\
 f_{\mathrm{heat}}(u)
 &=e^{-u},
 \\
 f_{\mathrm{Gauss}}(u)
 &=e^{-u^2},
 \\
 f_{\mathrm{heat2}}(u)
 &=(1+u)e^{-u}.
\end{align}
Их moments равны
\begin{align}
 \frac{f_2}{f_0}\bigg|_{\mathrm{sharp}}
 &=1,
 &
 \frac{f_2}{f_0}\bigg|_{\mathrm{heat}}
 &=1,
 \\
 \frac{f_2}{f_0}\bigg|_{\mathrm{Gauss}}
 &=\frac{\sqrt\pi}{2},
 &
 \frac{f_2}{f_0}\bigg|_{\mathrm{heat2}}
 &=2.
\end{align}

\section{Геометрический lock \(\Lambda R=1\)}

Как наиболее жёсткий no-new-scale test временно положим
\begin{equation}
 \Lambda R=1.
\end{equation}
Тогда
\begin{center}
\begin{tabular}{lll}
\toprule
Профиль & \((\chi R)^2\) & Vacuum status \\
\midrule
sharp & \(1/2\) & pass \\
heat & \(1/2\) & pass \\
Gauss & \((\sqrt\pi-1)/2\) & pass \\
heat2 & \(3/2\) & pass \\
\bottomrule
\end{tabular}
\end{center}
Все четыре естественные ветви дают nonzero vacuum, но три различных
значения \(\chi R\). Positivity, normalization и rapid decay не выбирают
единственную ветвь.

\begin{proposition}[Spectral-menu nonuniqueness]
Даже после геометрической фиксации \(\Lambda R=1\) стандартное
неотрицательное spectral-function menu не предсказывает уникальный
\(\chi R\). Совпадение sharp и heat profiles является частным равенством
их первых moments, а не общим selector theorem.
\end{proposition}

\section{Непрерывная profile--cutoff degeneracy}

Рассмотрим уже внутри heat-класса семейство
\begin{equation}
 f_a(u)=e^{-au},
 \qquad a>0.
\end{equation}
Оно удовлетворяет
\begin{equation}
 f_a(0)=1,
 \qquad
 \frac{f_2}{f_0}=\frac1a.
\end{equation}
Тогда
\begin{equation}
 \zeta_{\mathrm{mom}}
 =\frac{(\Lambda R)^2}{a}.
\end{equation}
Замена
\begin{equation}
 a\longmapsto s^2a,
 \qquad
 \Lambda R\longmapsto s\Lambda R
\end{equation}
оставляет \(\zeta_{\mathrm{mom}}\) неизменной. Следовательно, width
spectral profile и cutoff normalization не являются двумя физически
различимыми данными: они образуют одну непрерывную orbit.

\begin{theorem}[Moment--cutoff orbit no-go]
Условия \(f(0)=1\), \(f\ge0\), smoothness и rapid decay не фиксируют
\(f_2/f_0\). Более того, rescaling аргумента spectral function точно
компенсируется rescaling \(\Lambda\). Parent action видит только
\(\zeta_{\mathrm{mom}}\), которая остаётся одной непрерывной
scale-setting комбинацией.
\end{theorem}

\section{Vacuum thresholds}

Для profile с moment ratio \(\rho_f=f_2/f_0\) nonzero vacuum требует
\begin{equation}
 \Lambda R>\frac1{\sqrt{2\rho_f}}.
\end{equation}
Для меню:
\begin{align}
 \Lambda R&>\frac1{\sqrt2}
 &&\text{sharp/heat},
 \\
 \Lambda R&>\pi^{-1/4}
 &&\text{Gauss},
 \\
 \Lambda R&>\frac12
 &&\text{heat2}.
\end{align}
Эти inequalities являются genuine consistency predictions, но не
выбирают actual vacuum scale.

\section{Вердикт}

\begin{theorem}[Spectral-function selector no-go]
Минимальное естественное menu и непрерывный heat-profile class не выделяют
единственный spectral moment. Compact \(a_2/a_4\) relation остаётся
сильной boundary formula, но не превращается в absolute-scale prediction
без одной calibration величины \(\zeta_{\mathrm{mom}}\).
\end{theorem}

После задания одного train scale формула
\begin{equation}
 \zeta_{\mathrm{mom}}
 =(\chi R)^2+\frac12
\end{equation}
определяет требуемый spectral moment-cutoff invariant. Это корректная
реконструкция model parameter из train input, а не blind prediction.

\begin{protocol}[Следующий этап]
Scale-setting search в минимальном spectral-function class исчерпан.
Следует заморозить:
\begin{enumerate}
 \item один train scale, определяющий \(\zeta_{\mathrm{mom}}\);
 \item один representative spectral profile внутри rescaling orbit;
 \item dimensionless blind outputs, не использованные при выборе
 \(\zeta_{\mathrm{mom}}\).
\end{enumerate}
Дальнейший приоритет возвращается к physical representation/readout gate.
\end{protocol}

Машинный аудит сохранён в
\path{s2t_v3_spectral_function_moment_menu_results.json}.